\documentclass[11pt]{article}

\usepackage{amsmath,amssymb,amsthm,mathtools}
\usepackage{geometry}
\usepackage{hyperref}
\usepackage{enumitem}
\geometry{margin=1in}
\hypersetup{hidelinks}

\newtheorem{theorem}{Theorem}
\newtheorem{proposition}{Proposition}
\newtheorem{lemma}{Lemma}
\newtheorem{corollary}{Corollary}
\theoremstyle{definition}
\newtheorem{definition}{Definition}
\theoremstyle{remark}
\newtheorem{remark}{Remark}

\title{Temporal Bell Inequalities and Global Consistency in Modal Triplet Theory\\
\large A Conditional Sequential-Instrument Analysis}
\author{Peter Nero}
\date{July 2026}

\begin{document}
\maketitle

\begin{abstract}
We give a typed analysis of Leggett--Garg inequalities in Modal Triplet
Theory (MTT).  Three distinctions are decisive.  First, the laboratory
measurement times $t_1<t_2<t_3$ belong to a supplied physical clock on a
causal spacetime; they are not the auxiliary stabilization parameter used
to construct a fixed point.  Second, pairwise temporal correlations are
generated by declared channels and sequential measurement instruments,
including their state-update maps.  Third, a Leggett--Garg violation
excludes a single context-independent joint distribution for the measured
pair protocols; it does not identify projection alone as the cause.

We prove the standard bound $-3\leq K\leq1$ from the common-joint-distribution
contract.  We prove conditional no-retro-signaling for a causal sequential
instrument model and distinguish it from noninvasive measurability and
no-signaling in time.  A simple comparison shows why projection or
global-history language without dynamics, instruments, and a probability
source can support both Leggett--Garg-satisfying and violating realizations;
this is a scope check, not a new projection theorem.  Finally, an explicit
qubit, unitary, and L\"uders-instrument packet gives
$K=3/2$.  This is a standard quantum operational realization.  It becomes
an MTT result only if a selected MTT source descends to that same packet.
Thus global consistency remains a possible mechanism for temporal
contextuality, but the present paper is a conditional reconstruction and
completion contract, not a generic MTT prediction of temporal Bell
violation.
\end{abstract}

\section*{Revision note for this edition}
\addcontentsline{toc}{section}{Revision note for this edition}
\begin{description}
\item[Supersedes.] \emph{Temporal Bell Inequalities and Global Consistency in
Modal Triplet Theory}, version~2.
\item[Reason.] The version~2 release established the correct theorem boundary
but did not include the final theorem-ownership cleanup or enough guidance
for a reader to connect the probability model, pair protocols, and qubit
benchmark.
\item[Resolution.] Version~3 adds a plain-language protocol map and
result-by-result interpretation, and recasts the projection comparison and
MTT bridge as explanatory checks rather than duplicate formal theorems.
\item[Retained result.] The common-joint-distribution Leggett--Garg bound and
the conditional no-retro-signaling proposition are unchanged.
\item[Remaining boundary.] MTT does not yet derive the physical clock,
preparation state, sequential instruments, or arbitrary-apparatus
probability law required for a generic temporal Bell prediction.
\end{description}

\section*{How to read this paper}
\addcontentsline{toc}{section}{How to read this paper}
The paper has two layers.  The operational layer asks what an actual
three-time experiment must specify before a temporal correlation can be
calculated.  The MTT layer then asks whether one selected upper source emits
that entire operational packet.  Keeping the layers separate prevents a
geometric projection or a global-consistency slogan from being mistaken for
a probability law.

\paragraph{The central picture in plain language.}
Imagine three possible laboratory times but three separate experimental
runs.  One run measures at $t_1$ and $t_2$, another at $t_2$ and $t_3$, and
the third at $t_1$ and $t_3$.  Each first measurement can disturb the state
seen by the second.  The quantity $K$ therefore compares correlations from
three declared contexts; it is not automatically the expectation of three
values that existed together before measurement.

\paragraph{Argument map.}
Sections~1--2 separate physical clock time from the auxiliary parameter of a
fixed-point construction.  Section~3 defines the state, channels, instruments,
and pair probabilities.  Section~4 proves the algebraic bound under one
common joint distribution.  Section~5 distinguishes no-retro-signaling from
noninvasive measurability and from no-signaling in time.  Section~6 states
what global consistency can model and what additional MTT source rows are
still required.  Section~7 executes a concrete qubit example, after which the
final sections compare spatial Bell tests and state the completion and
falsification contract.

\tableofcontents

\section{Scope and Claim Tier}

Leggett--Garg inequalities test a conjunction of assumptions about a single
system measured at different physical times~\cite{LeggettGarg1985}.  In their
usual macrorealist interpretation, macrorealism per se is combined with
noninvasive measurability and an induction or causal-order assumption.
Experimental interpretation must also control the measurement-invasiveness
or ``clumsiness'' loophole~\cite{WildeMizel2012,EmaryLambertNori2014}.

The exact contribution of this paper is narrower:
\begin{enumerate}[label=\arabic*.]
\item identify the common-joint-distribution assumption that implies the
  three-time bound;
\item define the channels and instruments that generate the measured
  correlations;
\item state sufficient conditions for operational no-retro-signaling;
\item show what a selected MTT realization would have to emit;
\item exhibit the standard qubit benchmark.
\end{enumerate}

The paper does not derive a Born state, a measurement instrument, physical
time, a causal channel, or an objective history-selection law from projection
alone.  The current MTT shadow-bridge result likewise treats instruments as
selected updates rather than inverse or projection maps~\cite{NeroShadowBridge}.

\section{Two Parameters That Must Not Be Identified}

\begin{definition}[Physical measurement time]
Let $Y_4$ be a supplied time-oriented Lorentzian spacetime, or let a declared
laboratory clock provide an equivalent ordered parameter.  The events
\[
 t_1<t_2<t_3
\]
are physical measurement times on that clock.  They label preparation,
propagation, coupling, and readout events.
\end{definition}

\begin{definition}[Stabilization parameter]
An auxiliary parameter $\tau\geq0$ may label an iteration or semiflow
\[
 \frac{d u}{d\tau}=-\mathcal G(u)
\]
used to construct a fixed point, equilibrium, or coherent representative.
It is a mathematical evolution parameter on the construction space.
\end{definition}

\begin{remark}[Type separation]
No equality $t=\tau$ is assumed.  A theorem may use $\tau\to\infty$ to
construct data that are then assigned to a physical experiment at finite
clock times $t_i$.  Such a construction does not make stabilization time a
physical clock, a Lorentzian coordinate, or the shared circle phase.  This is
the same separation enforced in the corrected MTT kinematics program
\cite{NeroProgramA1}.
\end{remark}

\section{Sequential Operational Data}

\subsection{Channels and instruments}

This section is the operational engine of the paper.  A channel tells us how
the unmeasured system propagates.  An instrument contains more information:
it gives both the probability of each outcome and the conditional state left
behind.  Temporal correlations depend on that state update, so replacing an
instrument by a projector symbol or by an outcome label throws away data that
the experiment actually uses.

Let $\mathcal H$ be a complex Hilbert space and let $\mathfrak T(\mathcal H)$
be its trace-class operators.  A state is a positive
$\rho\in\mathfrak T(\mathcal H)$ with $\operatorname{Tr}\rho=1$.

\begin{definition}[Instrument]
For measurement time $t_i$ and setting $x_i$, a quantum instrument is a
family
\[
 \left\{\mathcal I^{(i)}_{q_i|x_i}:
 \mathfrak T(\mathcal H)\longrightarrow\mathfrak T(\mathcal H)
 \right\}_{q_i\in\{\pm1\}}
\]
of completely positive trace-nonincreasing maps such that the nonselective
map
\[
 \mathcal I^{(i)}_{x_i}:=
 \sum_{q_i=\pm1}\mathcal I^{(i)}_{q_i|x_i}
\]
is trace preserving.  Outcome probability and conditional state are
\[
 p(q_i|x_i)=\operatorname{Tr}\mathcal I^{(i)}_{q_i|x_i}(\rho),
 \qquad
 \rho_{q_i|x_i}
 =\frac{\mathcal I^{(i)}_{q_i|x_i}(\rho)}{p(q_i|x_i)}
\]
when the denominator is nonzero.
\end{definition}

This definition includes both an effect and an update law.  A POVM alone does
not determine the post-measurement state.  Instruments are standard
operational objects~\cite{DaviesLewis1970}.

Let $\mathcal E_{ji}$ be a completely positive trace-preserving channel from
time $t_i$ to $t_j$.  In a two-time protocol with common preparation
$\rho_i$, the pair probability is
\begin{equation}
p_{ij}(q_i,q_j|x_i,x_j)
=
\operatorname{Tr}\!\left[
\mathcal I^{(j)}_{q_j|x_j}\!\left(
\mathcal E_{ji}\!\left(
\mathcal I^{(i)}_{q_i|x_i}(\rho_i)
\right)\right)
\right].
\label{eq:pair}
\end{equation}
The measured correlator is
\begin{equation}
C_{ij}(x_i,x_j)
=
\sum_{q_i,q_j=\pm1}q_iq_j\,
p_{ij}(q_i,q_j|x_i,x_j).
\label{eq:correlator}
\end{equation}

\subsection{Pair protocols are contexts}

In the standard three-term test, $C_{12}$, $C_{23}$, and $C_{13}$ may be
estimated in distinct experimental runs.  In the $C_{13}$ run, the
measurement at $t_2$ is normally omitted.  Therefore the three pair
distributions are not automatically marginals of one observed three-time
distribution.  That common marginal structure is precisely part of the
classical contract tested by the inequality.

\section{The Leggett--Garg Bound}

The following theorem is deliberately algebraic.  It does not assume quantum
mechanics or MTT.  It identifies the exact classical data contract whose
failure is witnessed when the measured value exceeds the bound.

\begin{definition}[Common-joint-distribution contract]
Three measured pair distributions satisfy the common-joint-distribution
contract when there is one probability distribution
\[
 \mu(q_1,q_2,q_3),\qquad q_i\in\{\pm1\},
\]
whose $(1,2)$, $(2,3)$, and $(1,3)$ marginals equal the three corresponding
pair-protocol distributions.
\end{definition}

\begin{theorem}[Three-time Leggett--Garg bound]
\label{thm:lgi}
If the measured pair distributions satisfy the common-joint-distribution
contract, then
\[
 -3\leq K:=C_{12}+C_{23}-C_{13}\leq1.
\]
\end{theorem}

\begin{proof}
For every $(q_1,q_2,q_3)\in\{\pm1\}^3$, define
\[
 k(q_1,q_2,q_3)=q_1q_2+q_2q_3-q_1q_3.
\]
If $q_1=q_3$, then $k=2q_1q_2-1\in\{1,-3\}$.  If
$q_1=-q_3$, then $k=1$.  Hence $-3\leq k\leq1$ pointwise.
Expectation with respect to $\mu$ gives the claimed bound, because the
pairwise expectations of $\mu$ are the measured $C_{ij}$.
\end{proof}

Thus an observed value $K>1$ rules out the common-joint-distribution contract
for the three measured pair protocols.

\begin{remark}[What the violation does not decide]
This consequence does not by itself determine whether the failed physical
assumption is macrorealism, noninvasive measurability, context independence,
preparation equivalence, or another auxiliary premise.  That interpretation
requires the instruments and protocol controls.  Measurement disturbance is
therefore not a side issue; it is part of the theorem's empirical interface.
\end{remark}

\section{Causal Marginals and Measurement Disturbance}

The temporal direction makes two questions easy to confuse.  A later choice
should not alter a probability already recorded in the past; that is the
no-retro-signaling statement below.  By contrast, an earlier measurement may
change a later state and later statistics.  Ruling out that forward
disturbance is the stronger noninvasive-measurability or NSIT problem.

\subsection{No-retro-signaling}

\begin{definition}[Operational no-retro-signaling]
For $i<j$, operational no-retro-signaling holds when the earlier marginal
\[
 \sum_{q_j}p_{ij}(q_i,q_j|x_i,x_j)
\]
is independent of the later setting $x_j$.
\end{definition}

\begin{proposition}[Conditional no-retro-signaling]
\label{prop:nrs}
Assume Equation~\eqref{eq:pair}, a trace-preserving propagation channel
$\mathcal E_{ji}$, a complete future instrument
$\sum_{q_j}\mathcal I^{(j)}_{q_j|x_j}$ for every $x_j$, and no dependence of
the preparation, earlier instrument, or pre-$t_j$ channel on the future
setting.  Then operational no-retro-signaling holds.
\end{proposition}

\begin{proof}
Summing Equation~\eqref{eq:pair} over $q_j$ and using trace preservation gives
\begin{align*}
\sum_{q_j}p_{ij}(q_i,q_j|x_i,x_j)
&=\operatorname{Tr}\!\left[
\mathcal I^{(j)}_{x_j}\!\left(
\mathcal E_{ji}\!\left(
\mathcal I^{(i)}_{q_i|x_i}(\rho_i)
\right)\right)
\right]\\
&=\operatorname{Tr}\mathcal I^{(i)}_{q_i|x_i}(\rho_i),
\end{align*}
which is independent of $x_j$.
\end{proof}

\begin{remark}
Proposition~\ref{prop:nrs} is conditional on a causal operational model.  It
is not a theorem that every globally constrained history measure has
setting-independent earlier marginals.  If a future setting is allowed to
change the preparation or the measure assigned to earlier records, the proof
does not apply.
\end{remark}

\subsection{No-signaling in time is different}

No-signaling in time (NSIT) asks whether inserting an earlier nonselective
measurement leaves a later marginal unchanged:
\[
\sum_{q_i}p_{ij}(q_i,q_j|x_i,x_j)
=p_j(q_j|x_j).
\]
Unlike Proposition~\ref{prop:nrs}, NSIT generally fails for invasive
measurements because the earlier nonselective instrument can change the
state sent forward.  Operational no-retro-signaling, NSIT, and
macrorealist noninvasive measurability are therefore distinct predicates.
Temporal correlations also differ structurally from spacelike Bell
correlations because forward signaling and disturbance are physically
available~\cite{BrierleyEtAl2015}.

\section{What Global Consistency Can and Cannot Supply}

\subsection{A typed history model}

An MTT history realization may be described by:
\begin{itemize}
\item a selected upper history space $\Omega_x$ for each laboratory setting
  schedule $x$;
\item a selected probability state or measure $\nu_x$;
\item record maps $R_{i,x}:\Omega_x\to\{\pm1\}$;
\item a descent map from upper records to the operational packet
  $(\rho,\mathcal E,\mathcal I)$.
\end{itemize}
If different pair protocols are assigned setting-dependent history spaces or
measures, their pair distributions need not be marginals of one common
$\mu$.  This is a possible model of temporal contextuality.  It does not,
however, force a violation: many context-dependent families still satisfy
$K\leq1$.

\paragraph{Non-selection check.}
A projection map, a fixed-point construction, or the bare assertion that
histories are globally constrained does not determine the sign or magnitude
of a Leggett--Garg violation.  Keep the descriptive projection equal to the
identity.  One compatible realization assigns
$q_1=q_2=q_3=1$, giving $K=1$; another uses the qubit channels and instruments
of Section~\ref{sec:qubit}, giving $K=3/2$.  The difference lies in the state,
dynamics, instruments, and probability law, none of which is selected by the
projection.  This is an application of the general source-separation result,
not a new projection theorem owned by this paper.

\subsection{MTT completion contract}

For a selected MTT realization to claim the temporal-correlation result, its
source must supply:
\begin{enumerate}[label=(\roman*)]
\item a physical clock and causal order for $t_1<t_2<t_3$;
\item a normalized state or history measure;
\item channels and complete instruments whose descent gives
  Equation~\eqref{eq:pair};
\item common preparation and protocol conventions for all three pairs;
\item the causal independence assumptions of Proposition~\ref{prop:nrs};
\item pair correlators with $K>1$.
\end{enumerate}
With those rows supplied, Theorem~\ref{thm:lgi} excludes a common joint
distribution and Proposition~\ref{prop:nrs} gives operational
no-retro-signaling.  This is a completion contract assembled from the two
owned results, not a third theorem.

\begin{remark}[Current status]
The general MTT projection vocabulary
does not yet discharge its state, instrument, and source rows.  A selected
finite recorder may close a narrower apparatus contract without proving the
statement for arbitrary temporal experiments.
\end{remark}

\section{Explicit Qubit Benchmark}

This section is a worked example of the operational definitions, not a new
MTT prediction.  The example is useful because every ingredient is visible:
the state, the unitary propagator, the two-outcome instrument, the three pair
protocols, and the final value of $K$.
\label{sec:qubit}

Let
\[
\rho_\ast=\frac{\mathbf1}{2},\qquad
Q=\sigma_z,\qquad
P_q=\frac{\mathbf1+q\sigma_z}{2}.
\]
Use the L\"uders instrument
\[
\mathcal I_q(\rho)=P_q\rho P_q
\]
and unitary propagation
\[
\mathcal E_\delta(\rho)
=U_\delta\rho U_\delta^\dagger,\qquad
U_\delta=\exp\!\left(-\frac{i\omega\delta}{2}\sigma_x\right).
\]
The first outcome is uniform.  Conditional on outcome $q$, the probability of
the later outcome $r$ is
\[
p(r|q)=\operatorname{Tr}
\left(P_rU_\delta P_qU_\delta^\dagger\right)
=\frac{1+qr\cos(\omega\delta)}{2}.
\]
Therefore
\begin{equation}
\begin{aligned}
C(\delta)
&=\sum_{q,r=\pm1}qr\,p(q)p(r|q)\\
&=\cos(\omega\delta).
\end{aligned}
\label{eq:qubitcorr}
\end{equation}

Estimate $C_{12}$, $C_{23}$, and $C_{13}$ in separate pair protocols with
equal adjacent separation $\Delta$.  Equation~\eqref{eq:qubitcorr} gives
\[
C_{12}=C_{23}=\cos\theta,\qquad
C_{13}=\cos(2\theta),\qquad
\theta=\omega\Delta.
\]
Hence
\[
K=2\cos\theta-\cos(2\theta)
=1+2\cos\theta\bigl(1-\cos\theta\bigr).
\]
At $\theta=\pi/3$,
\[
K=\frac32.
\]

The nonselective L\"uders map leaves $\rho_\ast$ invariant, so the one-time
marginals in this benchmark are uniform.  The selective instrument still
updates the conditional state.  The packet is thus explicit enough to
evaluate both the correlations and the update law.

\begin{remark}[Tier of the benchmark]
This calculation is a standard quantum realization, not a derivation of
quantum instruments from MTT.  It is reproduced by an MTT model only when a
selected source and descent theorem emit the same state, channel, instrument,
clock convention, and correlators without importing the target values.
\end{remark}

\section{Relation to Spatial Bell Tests}

Spatial Bell and temporal Leggett--Garg inequalities both diagnose failures
of declared factorization contracts.  The analogy ends there unless a typed
intertwiner is supplied.  Spacelike Bell protocols impose two-way
no-signaling between separated laboratories.  Sequential temporal protocols
allow forward disturbance and communication through the system.  Therefore
one cannot transfer Jarrett's spatial decomposition unchanged to time, nor
infer a theorem of spatial--temporal unification from the word
``projection.''

Global consistency may provide a common interpretive language for the two
domains.  A physical unification theorem would additionally require common
source data, explicit maps to both operational scenarios, preservation of
their probability laws, and a proof that the relevant causal constraints
commute with both descents.

\section{Completion and Falsification Contract}

A physical MTT claim about temporal Bell experiments should report:
\begin{enumerate}[label=\textbf{T\arabic*.}]
\item the physical clock and causal channel;
\item the preparation state or normalized history measure;
\item every measurement instrument, including nonselective and conditional
  updates;
\item which pair protocols are run and which measurements are omitted;
\item the measured or predicted $C_{ij}$ and uncertainty budget;
\item no-retro-signaling and NSIT tests as separate rows;
\item an invasiveness-control protocol;
\item the MTT source path and descent certificate.
\end{enumerate}

The proposal fails at its declared tier if the source does not emit the
instrument packet, if protocol-dependent preparations are hidden, if
earlier marginals depend on later settings contrary to the claimed causal
contract, or if the predicted $K$ and auxiliary disturbance tests disagree
with held-out data.

\section{Previous Revision Record (Version 2)}

Version 2 makes the following corrections:
\begin{enumerate}[label=\arabic*.]
\item physical measurement times are separated from stabilization time;
\item measurement instruments and update rules are explicit;
\item the LGI theorem is stated for the common-joint-distribution contract;
\item violation is not attributed to projection or global consistency alone;
\item no-retro-signaling is conditional on a causal sequential model;
\item NSIT, no-retro-signaling, and noninvasive measurability are separated;
\item the qubit calculation is retained as an imported operational benchmark;
\item the former generic temporal-network and spatial--temporal unification
  claims are withdrawn.
\end{enumerate}

\section{Conclusion}

The robust result is the theorem boundary.  A common context-independent
joint distribution implies $K\leq1$.  Explicit sequential instruments are
needed to calculate the pair correlations and to identify disturbance.
Causal channels and complete future instruments imply operational
no-retro-signaling, but this does not establish noninvasive measurability.
The qubit L\"uders packet gives $K=3/2$ exactly.

MTT can host this packet and can use global-history data to model context
dependence.  It does not predict the violation from projection, fixed points,
or global consistency alone.  The remaining scientific step is a selected
same-source descent from MTT geometry and dynamics to the state, channels,
instruments, and uncertainty-controlled observables of a temporal experiment.

The practical lesson is correspondingly modest and useful.  Temporal Bell
language becomes testable only after every probability-bearing operation is
typed.  Once that is done, the ordinary inequality, the causal marginal
statement, and the qubit violation fit together without assigning any special
physical role to measurement itself.  Measurement is simply the declared
system--apparatus process that creates the record used in the correlation.

\begingroup
\raggedright
\begin{thebibliography}{99}

\bibitem{LeggettGarg1985}
A.~J.~Leggett and A.~Garg,
``Quantum mechanics versus macroscopic realism: Is the flux there when nobody
looks?,''
\emph{Physical Review Letters} \textbf{54}, 857--860 (1985),
\href{https://doi.org/10.1103/PhysRevLett.54.857}
{doi:10.1103/PhysRevLett.54.857}.

\bibitem{EmaryLambertNori2014}
C.~Emary, N.~Lambert, and F.~Nori,
``Leggett--Garg inequalities,''
\emph{Reports on Progress in Physics} \textbf{77}, 016001 (2014),
\href{https://doi.org/10.1088/0034-4885/77/1/016001}
{doi:10.1088/0034-4885/77/1/016001}.

\bibitem{WildeMizel2012}
M.~M.~Wilde and A.~Mizel,
``Addressing the clumsiness loophole in a Leggett--Garg test of
macrorealism,''
\emph{Foundations of Physics} \textbf{42}, 256--265 (2012),
\href{https://doi.org/10.1007/s10701-011-9598-4}
{doi:10.1007/s10701-011-9598-4}.

\bibitem{DaviesLewis1970}
E.~B.~Davies and J.~T.~Lewis,
``An operational approach to quantum probability,''
\emph{Communications in Mathematical Physics} \textbf{17}, 239--260 (1970).

\bibitem{BrierleyEtAl2015}
S.~Brierley, A.~Kosowski, M.~Markiewicz, T.~Paterek, and A.~Przysiezna,
``Nonclassicality of temporal correlations,''
\emph{Physical Review Letters} \textbf{115}, 120404 (2015),
\href{https://doi.org/10.1103/PhysRevLett.115.120404}
{doi:10.1103/PhysRevLett.115.120404}.

\bibitem{NeroShadowBridge}
P.~Nero,
``Projection, Probability, and Irreversibility: A Descent, Recovery, and
Measure-Separation Framework for MTT Shadow Bridges,'' version 3 (2026),
current selected MTT source.

\bibitem{NeroProgramA1}
P.~Nero,
``The Modal Triplet Theory Program A1: Coherent Kinematics,'' version 2
(2026), current selected MTT source.

\end{thebibliography}
\endgroup

\end{document}
